% GENERATED FILE. Edit canonical lessons, then run npm run generate:books.
% canonical-source-hash: fnv1a64:ac0fbf2c29e0669d
% canonical-lessons: AR-C93-shurta, AR-C93-suq, AR-C93-qarya, AR-C93-madina, AR-C93-matam

\chapter{Police, Market, Village, City, Restaurant}
\label{ch:ar-police-market-village-city-restaurant}
\hlchaptermodality{93}

\begin{chapteropening}
\textbf{By the end of this chapter:} \emph{I can say the police, a market, a village, a city and a restaurant.}

Complete the last lesson of chapter 93: I can say the police, a market, a village, a city and a restaurant.
\end{chapteropening}

\section[\texorpdfstring{\ar{شرطة} (shurṭa)}{شرطة (shurṭa)}]{\ar{شرطة} (shurṭa) — the police}

\label{lesson:AR-C93-shurta}



\begin{quote}
Before the new one: say the Arabic for a hotel, then the Arabic for an office; a desk.
\end{quote}

\subsection*{You'll want to know: \ar{شرطة}}
\textbf{\ar{شرطة}} — \emph{shurṭa} — \textquotedblleft{}the police\textquotedblright{}.

\begin{cousinweb}
\textbf{\ar{شرطة}} (\emph{shurṭa}) is built on the native Arabic root \emph{sh-r-ṭ}, \textquotedblleft{}to make a condition; a sign, a mark\textquotedblright{} (Lane). Lane's \emph{shurṭa} is the picked men of an army; the police are a later sense.
\end{cousinweb}

\begin{grammarlens}[title={What you've built}]
One more word for this chapter.
\end{grammarlens}

\subsection*{Your turn}
\begin{itemize}
\raggedright
  \item \emph{Say it:} \emph{shurṭa}
  \item \emph{Say it:} \emph{shurṭa}, once more
  \item \emph{Say it:} say \emph{shurṭa}
  \item \emph{Recall:} say the Arabic for a hotel, then the Arabic for an office; a desk, then say \emph{shurṭa} again
\end{itemize}

\begin{tcolorbox}[breakable,colback=teal!4,colframe=teal!35!black,title={Before you move on}]
What does \ar{شرطة} mean? (\textquotedblleft{}The police\textquotedblright{}.) Say it once more.
\end{tcolorbox}
\section[\texorpdfstring{\ar{سوق} (sūq)}{سوق (sūq)}]{\ar{سوق} (sūq) — a market}

\label{lesson:AR-C93-suq}



\begin{quote}
Before the new one: say the Arabic for an office; a desk, then the Arabic for the police.
\end{quote}

\subsection*{You'll want to know: \ar{سوق}}
\textbf{\ar{سوق}} — \emph{sūq} — \textquotedblleft{}a market\textquotedblright{}.

\begin{cousinweb}
\textbf{\ar{سوق}} (\emph{sūq}) is built on the native Arabic root \emph{s-w-q}, \textquotedblleft{}to drive (cattle)\textquotedblright{} (Lane). Lane quotes the explanation that a market is named because people drive their goods to it; English \emph{souk} comes from \emph{sūq}.
\end{cousinweb}

\begin{grammarlens}[title={What you've built}]
One more word for this chapter.
\end{grammarlens}

\subsection*{Your turn}
\begin{itemize}
\raggedright
  \item \emph{Say it:} \emph{sūq}
  \item \emph{Say it:} \emph{sūq}, once more
  \item \emph{Say it:} say \emph{sūq}
  \item \emph{Recall:} say the Arabic for an office; a desk, then the Arabic for the police, then say \emph{sūq} again
\end{itemize}

\begin{tcolorbox}[breakable,colback=teal!4,colframe=teal!35!black,title={Before you move on}]
What does \ar{سوق} mean? (\textquotedblleft{}A market\textquotedblright{}.) Say it once more.
\end{tcolorbox}
\section[\texorpdfstring{\ar{قرية} (qarya)}{قرية (qarya)}]{\ar{قرية} (qarya) — a village}

\label{lesson:AR-C93-qarya}



\begin{quote}
Before the new one: say the Arabic for the police, then the Arabic for a market.
\end{quote}

\subsection*{You'll want to know: \ar{قرية}}
\textbf{\ar{قرية}} — \emph{qarya} — \textquotedblleft{}a village\textquotedblright{}.

\begin{cousinweb}
\textbf{\ar{قرية}} (\emph{qarya}) is built on the native Arabic root \emph{q-r-y}, \textquotedblleft{}a town, a village\textquotedblright{} (Lane). Lane's \emph{qarya} is a town or village, smaller than a \emph{madīna}.
\end{cousinweb}

\begin{grammarlens}[title={What you've built}]
One more word for this chapter.
\end{grammarlens}

\subsection*{Your turn}
\begin{itemize}
\raggedright
  \item \emph{Say it:} \emph{qarya}
  \item \emph{Say it:} \emph{qarya}, once more
  \item \emph{Say it:} say \emph{qarya}
  \item \emph{Recall:} say the Arabic for the police, then the Arabic for a market, then say \emph{qarya} again
\end{itemize}

\begin{tcolorbox}[breakable,colback=teal!4,colframe=teal!35!black,title={Before you move on}]
What does \ar{قرية} mean? (\textquotedblleft{}A village\textquotedblright{}.) Say it once more.
\end{tcolorbox}
\section[\texorpdfstring{\ar{مدينة} (madīna)}{مدينة (madīna)}]{\ar{مدينة} (madīna) — a city}

\label{lesson:AR-C93-madina}



\begin{quote}
Before the new one: say the Arabic for a market, then the Arabic for a village.
\end{quote}

\subsection*{You'll want to know: \ar{مدينة}}
\textbf{\ar{مدينة}} — \emph{madīna} — \textquotedblleft{}a city\textquotedblright{}.

\begin{cousinweb}
\textbf{\ar{مدينة}} (\emph{madīna}) is listed by Lane under \emph{d-y-n}, \textquotedblleft{}to be subject, to be under authority\textquotedblright{}, as a city held under rule. Aramaic \emph{mdīnā}, \textquotedblleft{}province, city\textquotedblright{}, is the same formation, and many treat the Arabic word as borrowed from it.
\end{cousinweb}

\begin{grammarlens}[title={What you've built}]
One more word for this chapter.
\end{grammarlens}

\subsection*{Your turn}
\begin{itemize}
\raggedright
  \item \emph{Say it:} \emph{madīna}
  \item \emph{Say it:} \emph{madīna}, once more
  \item \emph{Say it:} say \emph{madīna}
  \item \emph{Recall:} say the Arabic for a market, then the Arabic for a village, then say \emph{madīna} again
\end{itemize}

\begin{tcolorbox}[breakable,colback=teal!4,colframe=teal!35!black,title={Before you move on}]
What does \ar{مدينة} mean? (\textquotedblleft{}A city\textquotedblright{}.) Say it once more.
\end{tcolorbox}
\section[\texorpdfstring{\ar{مطعم} (maṭʿam)}{مطعم (maṭʿam)}]{\ar{مطعم} (maṭʿam) — a restaurant}

\label{lesson:AR-C93-matam}



\begin{quote}
Before the new one: say the Arabic for a village, then the Arabic for a city.
\end{quote}

\subsection*{You'll want to know: \ar{مطعم}}
\textbf{\ar{مطعم}} — \emph{maṭʿam} — \textquotedblleft{}a restaurant\textquotedblright{}.

\begin{cousinweb}
\textbf{\ar{مطعم}} (\emph{maṭʿam}) is built on the native Arabic root \emph{ṭ-ʿ-m}, \textquotedblleft{}to eat, to taste\textquotedblright{} (Lane). \emph{Maṭʿam} is the \emph{ma-} place noun, \textquotedblleft{}a place of eating\textquotedblright{}.
\end{cousinweb}

\begin{grammarlens}[title={What you've built}]
One more word for this chapter.
\end{grammarlens}

\subsection*{Your turn}
\begin{itemize}
\raggedright
  \item \emph{Say it:} \emph{maṭʿam}
  \item \emph{Say it:} \emph{maṭʿam}, once more
  \item \emph{Say it:} say \emph{maṭʿam}
  \item \emph{Recall:} say the Arabic for a village, then the Arabic for a city, then say \emph{maṭʿam} again
\end{itemize}

\begin{tcolorbox}[breakable,colback=teal!4,colframe=teal!35!black,title={Before you move on}]
What does \ar{مطعم} mean? (\textquotedblleft{}A restaurant\textquotedblright{}.) Say it once more.
\end{tcolorbox}
